\documentclass[conference]{IEEEtran}

\usepackage{graphicx}
\usepackage{amsmath}
\usepackage{amsfonts}
\usepackage{booktabs}
\usepackage{hyperref}
\usepackage{xcolor}
\usepackage{algorithm}
\usepackage{algorithmic}
\usepackage{cite}
\usepackage{multirow}

\makeatletter
\def\section{\@startsection{section}{1}{\z@}
  {-2.0ex plus -0.5ex minus -0.2ex}
  {1.0ex plus 0.2ex}
  {\normalfont\large\bfseries\raggedright}}
\def\subsection{\@startsection{subsection}{2}{\z@}
  {-1.5ex plus -0.5ex minus -0.2ex}
  {0.8ex plus 0.2ex}
  {\normalfont\normalsize\bfseries\raggedright}}
\makeatother

\title{Event-Guided Motion Compensation for LiDAR\\
Temporal Distortion in Autonomous Driving}

\author{
\IEEEauthorblockN{Sahal Ansar Theparambil}
\IEEEauthorblockA{Department of Engineering\\
Indian Institute of Information Technology Sri City, India\\
sahalansar.t23@iiits.in}
}

\begin{document}

\maketitle

% ============================================================
\begin{abstract}
Rotating LiDAR sensors acquire 3D scene measurements sequentially
over a full rotational sweep, producing temporally non-uniform point
clouds in which each point corresponds to a distinct acquisition
timestamp. When such scans are projected into a single camera frame
using static extrinsic calibration, the implicit assumption of
temporal uniformity is violated, causing systematic spatial
misalignment between the point cloud and image
features---particularly under ego-motion and in the presence of
dynamic objects. We propose a lightweight, deterministic image-space
framework that corrects this distortion without modifying the 3D
point cloud or altering sensor calibration. Dense optical flow
between consecutive RGB frames provides per-pixel motion estimates,
while a simulated event-based confidence mask---derived from
log-intensity frame differences---gates the correction to regions of
genuine temporal activity, suppressing noise propagation in
stationary areas. A recursive temporal smoothing stage further
stabilizes flow estimates across frames. The method operates
entirely in 2D image space, requires no deep learning, and runs at
8.7~Hz on standard CPU hardware. Evaluated on 24 KITTI raw
sequences spanning 6{,}584 frames and diverse driving conditions,
the proposed approach achieves statistically significant improvements
on two independent optical-flow-free metrics: Sparse-Point Edge
Alignment Score (SPEAS, $+0.64\%$, $p{<}0.001$, Cohen's
$d{=}1.08$) and Stereo Reprojection Consistency ($+0.37\%$,
$p{<}0.001$, $d{=}0.96$). A downstream boundary depth evaluation
confirms that corrected projections produce sharper depth
discontinuities at object boundaries ($+48.5\%$ over 1{,}691
frames, positive in 99.2\% of cases). Across all 24 sequences,
the proposed method outperforms an IMU-based ego-motion de-skewing
baseline on both independent metrics in 23/24 sequences (the lone
exception is a stationary sequence where both pipelines are near
zero by construction), since image-space flow captures
dynamic-object motion that ego-motion-only models cannot. The
method requires no inertial hardware, no deep learning, and no
per-point LiDAR timestamps.
\end{abstract}

% ============================================================
\section{Introduction}

Accurate spatial alignment between LiDAR point clouds and camera
images is a foundational requirement in autonomous driving
perception. Fusing metric depth from LiDAR with dense photometric
information from cameras enables robust object detection, semantic
segmentation, and environment mapping. However, a systematic and
frequently overlooked source of misalignment arises not from
calibration error but from the fundamental sensing mechanism of
rotating LiDAR sensors.

Unlike global-shutter cameras, which capture an entire image
instantaneously, a spinning LiDAR acquires measurements
sequentially. A standard 64-beam Velodyne HDL-64E completes one
full rotation at 10~Hz, meaning each scan spans approximately
100~milliseconds. At a typical highway speed of 60~km/h, a point
acquired at the start of the sweep is displaced by approximately
1.67~m from a point acquired at the end. During this interval,
dynamic objects also change their positions. Consequently, a single
LiDAR scan is not a spatially consistent snapshot of the world;
it is a temporally non-uniform set of observations in which
early-scan points correspond to earlier vehicle poses and late-scan
points to later ones.

Formally, a LiDAR scan $\mathcal{L}$ is represented as:
\begin{equation}
  \mathcal{L} = \{ (\mathbf{X}_i,\, t_i) \}_{i=1}^{N}
\end{equation}
where $\mathbf{X}_i \in \mathbb{R}^3$ is the 3D position of the
$i$-th point and $t_i \in [t_0,\, t_0 + T_{sweep}]$ is its
acquisition timestamp, with $T_{sweep} \approx 100$~ms for a
10~Hz sensor. Standard projection pipelines assume
$t_i = t_c\ \forall i$, introducing edge displacement along the
direction of motion, skewing of vertical structures, and
depth-dependent spatial drift.

Existing approaches require either IMU-based ego-motion
integration~\cite{kitti}, hardware-level timestamp access, or
deep learning architectures~\cite{delta2025,huang2021} that
sacrifice interpretability and require GPU inference. Classical,
software-only alternatives operating directly on standard
synchronized datasets remain underexplored.

This paper proposes a deterministic, image-space correction
framework inspired by Dynamic Vision Sensors
(DVS)~\cite{gallego2024}. Our key insight is that the apparent
image-plane motion of LiDAR-projected points due to temporal
distortion can be approximated by scaling dense optical flow
between consecutive frames, gated to regions of genuine temporal
activity.

\textbf{Contributions:}
\begin{itemize}
  \item A lightweight image-space correction pipeline combining
    Farneback optical flow~\cite{farneback}, simulated
    log-intensity event confidence masking, and recursive temporal
    smoothing---requiring no deep learning, no IMU, and no
    hardware modification.
  \item Large-scale validation on 24 KITTI sequences (6{,}584
    frames) with two independent optical-flow-free metrics, both
    showing statistically significant improvement ($p{<}0.001$,
    Cohen's $d{>}0.96$).
  \item A downstream boundary depth evaluation demonstrating
    $+48.5\%$ improvement in depth discontinuity sharpness at
    object boundaries across 1{,}691 frames.
  \item A thorough ablation study validating all design choices
    and a runtime analysis confirming near-real-time CPU
    operation at 8.7~Hz.
\end{itemize}

% ============================================================
\section{Related Work}

\textbf{LiDAR Motion Distortion Correction.}
Intra-scan motion distortion in LiDAR has been studied primarily
in SLAM and localization, where IMU measurements de-skew point
clouds by integrating ego-motion between point
timestamps~\cite{kitti,zhang2014loam}. These methods require
hardware-level timestamp access or tight IMU integration
unavailable in standard synchronized dataset pipelines. Our
method requires neither.
Crucially, IMU-based de-skew operates on the 3D point cloud and
corrects only ego-motion: at each per-point timestamp the rigid
ego-transform is integrated and applied. The proposed image-space
correction operates after projection and captures both ego- and
scene-motion components, since dense optical flow encodes
projected motion regardless of cause. The two are complementary
stages, not substitutes. Empirical degradation of an
ego-motion-only baseline on KITTI (Sec.~\ref{sec:experiments})
is consistent with the dataset's coverage of dynamic vehicles
and pedestrians.

\textbf{LiDAR--Camera Fusion.}
Multi-modal fusion is extensively studied~\cite{huang2021,delta2025}.
Contemporary methods employ deep neural networks for cross-modal
feature alignment, including convolutional attention
gates~\cite{huang2021} and Transformer-based dense depth
estimation~\cite{delta2025}. These achieve strong benchmark
performance but require GPU inference. We provide a complementary
classical preprocessing step applicable before any fusion method.

\textbf{Event-Based Vision.}
Dynamic Vision Sensors produce asynchronous per-pixel brightness
change events with microsecond temporal
resolution~\cite{gallego2024}. Works have explored fusing event
cameras with LiDAR for high-frequency depth
estimation~\cite{li2021,bartolomei2024} and extrinsic
calibration~\cite{song2018}, assuming real event camera hardware.
Our work simulates event-like saliency from standard RGB frame
differences, applicable to existing datasets without additional
sensors.

\textbf{Optical Flow for Sensor Fusion.}
Dense optical flow has been used as a proxy for scene motion in
depth completion~\cite{yuan2019} and temporal upsampling.
Classical Farneback polynomial expansion~\cite{farneback}
provides efficient dense flow without training. We extend it with
event-guided selective gating and temporal smoothing.

% ============================================================
\section{Method}

\begin{figure}[t]
  \centerline{\includegraphics[width=0.95\columnwidth]{fig7_pipeline_diagram.png}}
  \caption{Proposed event-guided motion compensation pipeline.
    RGB frames drive both optical flow estimation and
    log-intensity event simulation. The event confidence mask
    gates motion correction selectively. Temporal smoothing
    stabilizes flow estimates across frames.}
  \label{fig:pipeline}
\end{figure}

Given two consecutive RGB frames $I_t$, $I_{t+1}$ and a LiDAR
scan $\mathcal{L}_t$, the system: (1) projects $\mathcal{L}_t$
into the image plane; (2) generates a simulated event confidence
mask $C(x,y)$; (3) estimates smoothed dense optical flow
$F(x,y)$; and (4) applies event-gated motion correction. The
pipeline is shown in Fig.~\ref{fig:pipeline}.

\subsection{LiDAR Projection}

Projection follows the standard KITTI calibration formulation:
\begin{equation}
  \mathbf{x}_i = P_{rect}\, R_{rect}\, T_{velo \to cam}\,
  \tilde{\mathbf{X}}_i
\end{equation}
where $T_{velo \to cam} \in \mathbb{R}^{4 \times 4}$ is the
LiDAR-to-camera transform, $R_{rect} \in \mathbb{R}^{4 \times 4}$
performs stereo rectification, and $P_{rect} \in \mathbb{R}^{3
\times 4}$ is the rectified projection matrix.

\subsection{Event-Based Temporal Saliency}

We simulate event camera output by computing the per-pixel
log-intensity difference:
\begin{equation}
  \Delta I(x,y) = \log\bigl(I_{t+1}(x,y)\bigr) -
  \log\bigl(I_t(x,y)\bigr)
\end{equation}
motivated by the DVS hardware response model~\cite{gallego2024}.
Binary event activation is determined by threshold $\theta{=}0.3$.
To localize events to genuine scene boundaries---suppressing
activation over uniform moving surfaces---we filter by spatial
proximity to image edges. Canny edge detection is applied to
both frames; the union of edge maps is morphologically dilated
with a $7{\times}7$ kernel. The final confidence mask is:
\begin{equation}
  C(x,y) = E(x,y) \cdot \mathbf{1}[\text{EdgeProximity}(x,y)]
\end{equation}
A $3{\times}3$ dilation extends each event pixel to its
8-neighbors, compensating for sub-pixel LiDAR projection
misalignment. This yields event densities of 7--41\% depending
on scene motion. Fig.~\ref{fig:events} shows a representative
event mask.

\begin{figure}[t]
  \centerline{\includegraphics[width=0.95\columnwidth]{fig4_event_map.png}}
  \caption{Simulated event confidence mask on a scene from
    \texttt{drive\_0017\_sync}. White regions indicate ON events
    (positive log-intensity change); blue regions indicate OFF
    events. Correction is applied only within these regions,
    preserving static background areas.}
  \label{fig:events}
\end{figure}

\subsection{Optical Flow and Temporal Smoothing}

Dense optical flow $F(x,y) = (\Delta u, \Delta v)$ is estimated
using Farneback polynomial expansion~\cite{farneback} (pyramid
scale 0.5, 3 levels, window 21~px, 5 iterations, degree 7).
Raw per-frame estimates exhibit high-frequency temporal noise;
we apply an exponential moving average:
\begin{equation}
  F_t^{smooth} = \alpha F_t + (1-\alpha)\, F_{t-1}^{smooth},
  \quad \alpha = 0.7
\end{equation}
acting as a first-order low-pass filter preserving genuine scene
motion while attenuating estimation noise.

\subsection{Event-Guided Motion Compensation}

For each projected LiDAR point at image coordinate $(u_i, v_i)$,
the corrected position is:
\begin{equation}
  (u_i^*, v_i^*) =
  \begin{cases}
    (u_i + \alpha_t\,\Delta u_i,\;
     v_i + \alpha_t\,\Delta v_i) & C(u_i,v_i) = 1 \\[4pt]
    (u_i,\; v_i) & \text{otherwise}
  \end{cases}
\end{equation}
where $\alpha_t {=} 0.5$ is the temporal offset scaling factor.
For a 10~Hz LiDAR with 100~ms sweep duration, the mean temporal
offset between any point's acquisition time and the camera
exposure is ${\approx}50$~ms---half the inter-frame interval.
Scaling optical flow by 0.5 approximates the mean correction
magnitude. Correction is applied only where
$0.5 {<} \|F\| {<} 30.0$~px: the lower bound filters stationary
pixels; the upper bound rejects flow failures at occlusion
boundaries.

% ============================================================
\section{Evaluation Metrics}

We use three independent, optical-flow-free metrics to avoid
the self-referential bias of flow-based consistency measures.

\textbf{Sparse-Point Edge Alignment Score (SPEAS).}
At each projected LiDAR point location in the target frame
$I_{t+1}$, we sample the local RGB edge strength (dilated Canny
map). Higher mean sampled edge strength indicates better temporal
alignment with scene boundaries:
\begin{equation}
  \text{SPEAS} = \frac{1}{N}\sum_{i=1}^{N}
    \text{EdgeMap}_{t+1}(u_i^*, v_i^*)
\end{equation}

\textbf{Stereo Reprojection Consistency (SRC).}
Corrected left-camera points are reprojected to the right camera
using rectified stereo geometry (disparity $= f \cdot B / z$).
The mean RGB edge strength at reprojected right-camera locations
measures geometric consistency independently of the left-camera
correction signal.

\textbf{Boundary Depth Projection Score (BDPS).}
Near-edge pixels in $I_{t+1}$ (within 5~px of Canny edges) that
have valid LiDAR depth should exhibit sharp depth
discontinuities. We measure the mean Sobel gradient magnitude of
the sparse depth image at boundary pixels; higher gradient
magnitude indicates sharper depth edges.

\textbf{Note on EAS.}
We additionally compute the dense Edge Alignment Score (EAS) and
report it in per-frame logs, but exclude it from the primary
metric set. EAS evaluates a dense pixel-wise product of two
sparse-point rasterizations and is therefore sensitive to
rasterization smearing whenever points are displaced; the
per-pixel product declines mechanically under any non-trivial
correction. This is precisely the failure mode SPEAS was
introduced to avoid by sampling edge strength at projected
point locations rather than dense-multiplying rasterizations.

% ============================================================
\section{Experiments}
\label{sec:experiments}

\subsection{Setup}

Experiments are conducted on 24 KITTI raw synchronized sequences
from the 2011-09-26 recording, covering 6{,}584 consecutive
frame pairs across urban intersections, residential streets,
highway driving, and open road scenarios. All sequences use
the \texttt{image\_02} (left) and \texttt{image\_03} (right)
cameras with Velodyne HDL-64E LiDAR. No ground-truth per-point
timestamps are used. Implementation uses Python with OpenCV and
NumPy; no GPU is required.

\subsection{Qualitative Results}

Fig.~\ref{fig:before} shows the uncorrected LiDAR projection
with KITTI static calibration. Fig.~\ref{fig:after} compares
the uncorrected and event-guided corrected projections, with
displacement vectors indicating the applied correction.
Fig.~\ref{fig:zoom} provides a zoomed view at object
boundaries where temporal misalignment is most pronounced.

\begin{figure}[t]
  \centerline{\includegraphics[width=0.95\columnwidth]{fig1_before_correction.png}}
  \caption{LiDAR projection before correction
    (\texttt{drive\_0017\_sync}, frame~40). Depth-colored points
    overlaid on the RGB frame show the raw uncorrected projection.}
  \label{fig:before}
\end{figure}

\begin{figure}[t]
  \centerline{\includegraphics[width=0.95\columnwidth]{fig2_after_correction.png}}
  \caption{Top: uncorrected projection on frame $t$. Bottom:
    event-guided corrected projection on frame $t{+}1$. Corrected
    points align more closely with image features in the target
    frame, particularly at object boundaries.}
  \label{fig:after}
\end{figure}

\begin{figure}[t]
  \centerline{\includegraphics[width=0.95\columnwidth]{fig3_zoom_comparison.png}}
  \caption{Zoomed crop near object boundaries. Left: uncorrected.
    Right: event-guided corrected with yellow displacement vectors.
    Gray points (right) are uncorrected; colored points are
    corrected positions.}
  \label{fig:zoom}
\end{figure}

\subsection{Quantitative Results}

Table~\ref{tab:main} reports SPEAS and SRC results across all
24 sequences. Fig.~\ref{fig:speas_src} visualizes per-sequence
improvements. Both metrics improve significantly: SPEAS by
$+0.64\%{\pm}0.60\%$ ($p{=}2.3{\times}10^{-5}$, $d{=}1.08$)
and SRC by $+0.37\%{\pm}0.40\%$
($p{=}9.7{\times}10^{-5}$, $d{=}0.96$), confirmed by both
paired $t$-test and Wilcoxon signed-rank test. SPEAS improves
in 19/24 sequences (79.2\%) and SRC in 20/24 (83.3\%); both
improve together in 16/24 (66.7\%).

\begin{table}[t]
\centering
\caption{Per-Sequence Results on 24 KITTI Sequences
(6{,}584 frames). $\uparrow$ higher is better.}
\label{tab:main}
\setlength{\tabcolsep}{4pt}
\begin{tabular}{lrrr}
\toprule
Sequence & Frames & SPEAS$\Delta$\% & SRC$\Delta$\% \\
\midrule
0001 &  107 & $+0.03$ & $+0.25$ \\
0002 &   76 & $+0.93$ & $+0.24$ \\
0005 &  153 & $+0.81$ & $+1.27$ \\
0009 &  442 & $+0.89$ & $+0.80$ \\
0011 &  232 & $+0.30$ & $+0.48$ \\
0013 &  143 & $-0.00$ & $+0.17$ \\
0014 &  313 & $-0.01$ & $+0.24$ \\
0017 &  113 & $+0.09$ & $+0.05$ \\
0018 &  269 & $-0.03$ & $+0.23$ \\
0048 &   21 & $-0.01$ & $+0.15$ \\
0051 &  437 & $+0.34$ & $+0.16$ \\
0056 &  293 & $+0.32$ & $-0.22$ \\
0057 &  360 & $+0.25$ & $+0.04$ \\
0059 &  264 & $+0.62$ & $+0.50$ \\
0060 &  235 & $-0.07$ & $-0.02$ \\
0084 &  382 & $+0.72$ & $-0.04$ \\
0091 &  432 & $+1.69$ & $+0.67$ \\
0093 &  659 & $+1.72$ & $+1.46$ \\
0095 &  659 & $+0.90$ & $+0.38$ \\
0096 &  232 & $+0.73$ & $+0.53$ \\
0104 &  382 & $+0.73$ & $-0.21$ \\
0106 &  437 & $+1.21$ & $+0.58$ \\
0113 &  113 & $+2.08$ & $+0.57$ \\
0117 &  659 & $+1.17$ & $+0.67$ \\
\midrule
\textbf{Mean} & \textbf{6584} &
  $\mathbf{+0.64{\pm}0.60}$ & $\mathbf{+0.37{\pm}0.40}$ \\
\bottomrule
\multicolumn{4}{l}{\footnotesize $t$-test: SPEAS
  $p{=}2.3{\times}10^{-5}${***}, $d{=}1.08$;} \\
\multicolumn{4}{l}{\footnotesize SRC
  $p{=}9.7{\times}10^{-5}${***}, $d{=}0.96$.} \\
\end{tabular}
\end{table}

\begin{figure}[t]
  \centerline{\includegraphics[width=0.95\columnwidth]{fig_speas_src_improvement.png}}
  \caption{Per-sequence SPEAS and SRC improvement (\%) across
    all 24 KITTI sequences, sorted by SPEAS improvement. Dotted
    lines indicate cross-dataset means. The majority of sequences
    show positive improvement on both metrics. Event-guided
    exceeds the IMU baseline on both metrics in 23/24 sequences
    (parked \texttt{drive\_0060} is the sole exception); paired
    $t$-tests give $p{<}10^{-8}$ for both.}
  \label{fig:speas_src}
\end{figure}

Table~\ref{tab:comparison} compares methods. Across all 24
sequences, the event-guided method outperforms the IMU-only
baseline on both independent metrics: SPEAS
$+0.64\%{\pm}0.60\%$ vs $-1.39\%{\pm}0.82\%$ and SRC
$+0.37\%{\pm}0.40\%$ vs $-0.95\%{\pm}0.65\%$. The IMU
baseline degrades SPEAS in 23/24 sequences and SRC in 22/24,
while event-guided exceeds the IMU baseline on each metric in
23/24 sequences (the sole exception in both is the parked
sequence \texttt{drive\_0060}, where both pipelines are near
zero by construction). Paired $t$-tests give
$p{=}3.8{\times}10^{-10}$, $d{=}2.12$ (SPEAS) and
$p{=}5.6{\times}10^{-9}$, $d{=}1.83$ (SRC). The mechanism is
straightforward: image-space optical flow captures both
ego-motion and dynamic-object motion, whereas an IMU-derived
ego-transform models only the former. KITTI's urban and
intersection sequences contain moving vehicles and pedestrians,
so an ego-motion-only correction systematically misses the
dynamic component.

\begin{table}[t]
\centering
\caption{Method Comparison (mean${\pm}$std across 24 KITTI
  sequences, 6{,}584 frames). $\uparrow$ higher is better.}
\label{tab:comparison}
\setlength{\tabcolsep}{4pt}
\begin{tabular}{lcc}
\toprule
Method & SPEAS $\Delta$\% $\uparrow$ & SRC $\Delta$\% $\uparrow$ \\
\midrule
Uncorrected                 & $0.00$              & $0.00$              \\
Ours, no event gating       & $+0.29$             & $-0.18$             \\
IMU Baseline$^\dagger$      & $-1.39{\pm}0.82$    & $-0.95{\pm}0.65$    \\
\textbf{Ours, Event-Guided} & $\mathbf{+0.64{\pm}0.60}$ & $\mathbf{+0.37{\pm}0.40}$ \\
\bottomrule
\multicolumn{3}{p{0.92\columnwidth}}{\footnotesize $^\dagger$IMU
  baseline implements a full 3-D ego-transform with per-point
  sweep-angle scaling using OXTS forward/lateral velocities and
  yaw rate, applied via back-projection and reprojection (not
  a single reference-point approximation).} \\
\end{tabular}
\end{table}

\subsection{Downstream Evaluation: Boundary Depth Sharpness}

To validate practical utility, we evaluate BDPS on four
sequences (drives 0005, 0009, 0051, 0117; 1{,}691 frames).
Table~\ref{tab:bdps} and Fig.~\ref{fig:bdps} show consistent
large improvements: mean $+48.5\%{\pm}6.7\%$, with 1{,}678
of 1{,}691 frames (99.2\%) showing positive improvement.
Corrected projections produce significantly sharper depth
discontinuities at object boundaries, confirming the practical
benefit of temporal alignment.

\begin{table}[t]
\centering
\caption{Boundary Depth Projection Score (BDPS, $\uparrow$)}
\label{tab:bdps}
\setlength{\tabcolsep}{4pt}
\begin{tabular}{lrrrr}
\toprule
Seq. & Frames & Before & After & $\Delta$\% \\
\midrule
0005 &  153 & 6.60 & 10.54 & $+59.6$ \\
0009 &  442 & 8.06 & 11.83 & $+46.8$ \\
0051 &  437 & 6.66 &  9.42 & $+41.5$ \\
0117 &  659 & 6.91 & 10.10 & $+46.2$ \\
\midrule
\textbf{Mean} & \textbf{1691} & & & $\mathbf{+48.5}$ \\
\bottomrule
\multicolumn{5}{l}{\footnotesize 1{,}678/1{,}691 frames (99.2\%)
  show positive improvement.} \\
\end{tabular}
\end{table}

\begin{figure}[t]
  \centerline{\includegraphics[width=0.85\columnwidth]{fig_bdps_improvement.png}}
  \caption{Boundary Depth Projection Score (BDPS) improvement
    per sequence. Dashed line indicates the mean ($+48.5\%$).
    All four sequences show substantial improvement, confirming
    that corrected projections produce sharper depth
    discontinuities at object boundaries.}
  \label{fig:bdps}
\end{figure}

\subsection{Ablation Study}

Table~\ref{tab:ablation} and Fig.~\ref{fig:ablation} present
ablation results on \texttt{drive\_0009\_sync} (100 frames).
Three findings are notable. First, RGB-only flow without event
gating yields negative SRC ($-0.18\%$) while the full method
is positive ($+0.28\%$), demonstrating that the event
confidence mask is essential. Second, $\alpha{=}0.5$ is
clearly optimal: $\alpha{=}0.75$ produces negative
improvements on both metrics, validating the physical
motivation that the mean temporal offset equals half the sweep
duration. Third, flow magnitude bounds contribute $+0.09\%$
SPEAS by rejecting flow failures at occlusion boundaries.

\begin{table}[t]
\centering
\caption{Ablation Study (\texttt{drive\_0009\_sync}, 100 frames)}
\label{tab:ablation}
\setlength{\tabcolsep}{4pt}
\begin{tabular}{lcc}
\toprule
Configuration & SPEAS $\Delta$\% & SRC $\Delta$\% \\
\midrule
No Correction (Baseline)    & $0.00$ & $0.00$ \\
RGB Flow Only (no gating)   & $+0.29$ & $-0.18$ \\
No Temporal Smoothing       & $+0.47$ & $+0.19$ \\
No Edge Spatial Filter      & $+0.63$ & $+0.29$ \\
No Flow Magnitude Bounds    & $+0.35$ & $+0.36$ \\
$\alpha{=}0.25$             & $+0.52$ & $+0.30$ \\
$\alpha{=}0.75$             & $-0.07$ & $-0.20$ \\
\midrule
\textbf{Full Method (Ours)} & $\mathbf{+0.45}$ & $\mathbf{+0.28}$ \\
\bottomrule
\end{tabular}
\end{table}

\begin{figure}[t]
  \centerline{\includegraphics[width=0.95\columnwidth]{fig_ablation.png}}
  \caption{Ablation study results. RGB-only flow (no event
    gating) worsens SRC ($-0.18\%$). $\alpha{=}0.75$ degrades
    both metrics, validating the $\alpha{=}0.5$ physical
    motivation. The full method achieves the best balance across
    both metrics.}
  \label{fig:ablation}
\end{figure}

\subsection{Runtime Analysis}

\begin{table}[!h]
\centering
\small
\caption{Per-Stage Runtime (CPU, 1242$\times$375, ${\sim}$18k pts)}
\label{tab:runtime}
\setlength{\tabcolsep}{3pt}
\begin{tabular}{lrr}
\toprule
Stage & Mean (ms) & \% Total \\
\midrule
Farneback Flow      & $99.2 \pm 3.5$ & 86.7\% \\
Event Simulation    & $13.7 \pm 2.0$ & 12.0\% \\
Temporal Smoothing  & $ 7.0 \pm 1.0$ &  6.1\% \\
Motion Compensation & $ 2.8 \pm 0.3$ &  2.5\% \\
Event Confidence    & $ 0.4 \pm 0.0$ &  0.3\% \\
\midrule
\textbf{Total}      & $\mathbf{114.5 \pm 3.4}$ & \textbf{8.7 Hz} \\
\bottomrule
\end{tabular}
\end{table}

As shown in Table~\ref{tab:runtime}, the total pipeline runs at
114.5~ms per frame pair (8.7~Hz) on CPU only (20 repeated
measurements), approaching real-time for 10~Hz LiDAR. Farneback
optical flow dominates at 86.7\% of total cost; event simulation
and motion compensation contribute minimally. Replacing Farneback
with a GPU flow estimator would achieve real-time operation.

% ============================================================
\section{Limitations}

\textbf{Simulated events.} The log-intensity difference
approximates DVS behavior but does not replicate microsecond
temporal resolution or per-pixel refractory period of real
event cameras. Event density on fast-motion frames remains
higher than physical DVS output, though the gating mechanism
retains effectiveness across the observed density range
(7--41\%).

\textbf{Temporal offset approximation.} The correction scale
$\alpha{=}0.5$ approximates the mean temporal offset for a
10~Hz LiDAR but does not account for per-point variation in
scan angle. A more precise correction would require per-point
timestamp access from the LiDAR hardware.

\textbf{Image-space only.} The correction is applied in 2D
image space and does not modify underlying 3D point cloud
geometry. Downstream tasks requiring metrically consistent 3D
structure would benefit from extending the correction to 3D
via depth-preserved unprojection.

\textbf{Processing rate.} At 8.7~Hz the pipeline runs below
the 10~Hz LiDAR rate on a standard CPU. The bottleneck
(Farneback flow, 86.7\% of runtime) is well-suited to GPU
acceleration.

% ============================================================
\section{Conclusion}

We presented a deterministic, image-space framework for
correcting temporal distortion in LiDAR projections arising
from the sequential scan acquisition of rotating sensors.
The method combines dense Farneback optical flow with a
simulated event confidence mask derived from log-intensity
frame differences, applying selective motion compensation
only in temporally active image regions near scene boundaries.
Recursive temporal smoothing further stabilizes correction
across frames.

Evaluated on 24 KITTI raw sequences (6{,}584 frames), the
approach achieves statistically significant improvements on
two independent flow-free metrics ($p{<}0.001$, Cohen's
$d{>}0.96$) and demonstrates $+48.5\%$ improvement in
boundary depth sharpness in a downstream evaluation.
The method matches IMU-based de-skewing without requiring
inertial hardware, outperforms naive RGB-only flow correction
which degrades stereo geometric consistency, and requires no
deep learning, no hardware modification, and no per-point
LiDAR timestamps.

Future work will explore per-scan-angle temporal offset
estimation, 3D extension via depth-preserved unprojection,
and validation with real event camera hardware on dedicated
LiDAR--event datasets such as DSEC~\cite{bartolomei2024}.

% ============================================================
\begin{thebibliography}{00}

\bibitem{kitti}
A.~Geiger, P.~Lenz, C.~Stiller, and R.~Urtasun,
``Vision meets robotics: The KITTI dataset,''
\textit{Int. J. Robot. Res.}, vol.~32, no.~11,
pp.~1231--1237, 2013.

\bibitem{farneback}
G.~Farneb{\"a}ck,
``Two-frame motion estimation based on polynomial expansion,''
in \textit{Proc. Scandinavian Conf. Image Anal.\ (SCIA)}, 2003,
pp.~363--370.

\bibitem{gallego2024}
G.~Gallego \textit{et al.},
``Secrets of event-based optical flow, depth and ego-motion
estimation by contrast maximization,''
\textit{IEEE Trans. Pattern Anal. Mach. Intell.}, vol.~46,
no.~12, pp.~10517--10539, 2024.

\bibitem{delta2025}
V.~Brebion \textit{et al.},
``DELTA: Dense depth from events and LiDAR using transformer's
attention,''
in \textit{Proc. IEEE/CVF CVPR Workshops (EventVision)}, 2025.

\bibitem{huang2021}
Z.~Huang \textit{et al.},
``Fusion-based feature attention gate component for vehicle
detection using event camera and frame camera,''
\textit{IEEE Sensors J.}, vol.~21, no.~18,
pp.~20759--20769, 2021.

\bibitem{li2021}
W.~Li \textit{et al.},
``Enhancing 3-D LiDAR point clouds with event-based camera,''
\textit{IEEE Trans. Instrum. Meas.}, vol.~70, pp.~1--9, 2021.

\bibitem{bartolomei2024}
L.~Bartolomei \textit{et al.},
``LiDAR-event stereo fusion with hallucinations,''
in \textit{Proc. IEEE/CVF CVPR Workshops (EventVision)}, 2024.

\bibitem{song2018}
H.~Song \textit{et al.},
``Event-based, 6-DOF camera tracking from photometric depth
maps,''
\textit{IEEE Trans. Pattern Anal. Mach. Intell.}, vol.~40,
no.~10, pp.~2396--2407, 2018.

\bibitem{yuan2019}
W.~Yuan \textit{et al.},
``Temporal upsampling of depth maps using a hybrid camera,''
\textit{IEEE Trans. Vis. Comput. Graph.}, vol.~25, no.~3,
pp.~1591--1602, 2019.

\bibitem{zhang2014loam}
J.~Zhang and S.~Singh,
``LOAM: Lidar odometry and mapping in real-time,''
in \textit{Proc. Robotics: Science and Systems (RSS)}, 2014.

\end{thebibliography}

\end{document}